\subsection{Большие языковые модели и инференс}

Большая языковая модель --- это нейронная модель, обученная на больших корпусах текстовых и близких к ним символьных данных и предназначенная для обработки последовательностей. Современные большие языковые модели основываются на архитектуре трансформера, предложенной в работе ``Attention Is All You Need''~\cite{vaswani2017attention}. В отличие от рекуррентных архитектур, трансформер использует механизм внимания, что позволяет эффективно учитывать зависимости между токенами и хорошо распараллеливать вычисления на GPU.

Рост качества больших языковых моделей тесно связан с масштабированием числа параметров, объема обучающих данных и вычислительного бюджета. На примере GPT-3 было показано, что увеличение масштаба авторегрессионной языковой модели приводит к появлению новых возможностей и повышению качества на широком наборе задач~\cite{brown2020language}. Работы о формализации масштабировании моделей и вычислительно оптимальном обучении показывают, что качество языковой модели зависит не только от размера модели, но и от соотнешения между числом ее параметров и объемом данных~\cite{kaplan2020scaling,hoffmann2022training}. Это объясняет, почему развитие больших языковых моделей постоянно упирается в системные ограничения: память GPU, пропускную способность между устройствами и стоимость выполнения.

Жизненный цикл большой языковой модели делится на два этапа: обучение и вывод. На этапе обучения параметры модели обновляются на основе больших наборов данных. Этот этап требует огромного числа прямых и обратных проходов, хранения градиентов и промежуточных активаций. Поэтому для обучения особенно важны масштабирование на многие GPU, быстрый обмен градиентами и возможность долго выполнять вычисления без отказов.

Этап вывода, или инференса, имеет другую природу. Параметры модели уже зафиксированы, а система должна отвечать на пользовательские запросы с ограничениями по задержке, пропускной способности и стабильности. Для авторегрессионных языковых моделей инференс выполняется последовательно по токенам: модель получает текущий контекст, вычисляет распределение следующего токена, добавляет выбранный токен в контекст и повторяет процесс. Поэтому даже небольшая задержка одного шага умножается на число токенов в выходной последовательности. На практике важны не только средняя скорость, но и верхние перцентили задержки, так как пользовательский запрос завершается только после выполнения всей цепочки шагов.

Авторегрессионная природа инференса делает один шаг коротким, но многократно повторяемым. На каждом шаге запускаются операции на GPU, выполняются обращения к KV кэшу, а в распределенной системе еще и передаются промежуточные тензоры между устройствами. Поэтому существенной частью задержки становится не только арифметика внутри слоев, но и накладные расходы запуска операций и согласования работы нескольких GPU.

Инференс крупных моделей осложняется тем, что веса модели, KV кэш при длинных контекстах и промежуточные активации при большом числе параллельных запросов могут не помещаться в память одного GPU. Даже если модель помещается на одно устройство, один GPU может не обеспечить требуемую пропускную способность при большом числе параллельных запросов. Поэтому промышленные системы инференса используют несколько устройств и совмещают вычислительные оптимизации, управление памятью и распределенное выполнение. Например, DeepSpeed Inference рассматривает инференс трансформеров как системную задачу, в которой нужно одновременно учитывать размер модели, задержку, пропускную способность, оптимизированные ядра и размещение данных между GPU, CPU и накопителями~\cite{aminabadi2022deepspeed}.

\subsection{Архитектура Mixture-of-Experts}

Один из способов масштабировать языковые модели --- увеличивать плотную модель, в которой каждый вход проходит через все имеющиеся параметры. Такой подход прост с точки зрения архитектуры, но дорог при инференсе: каждый прямой проход использует весь набор весов. Поэтому увеличение плотной модели почти напрямую повышает стоимость одного шага инференса.

Архитектура Mixture-of-Experts, или MoE, предлагает другой путь. В ней полносвязные линейные слои (feed-forward network) в трансформере заменяются наборами экспертов, а отдельный роутер выбирает, какие эксперты должны обработать конкретный токен или группу токенов. В работе Shazeer et al. был предложен sparsely-gated MoE слой, где для каждого входа активируется только небольшое число экспертов~\cite{shazeer2017moe}. За счет этого можно увеличить общее число параметров модели, но не использовать все экспертные параметры при каждом проходе.

Главное преимущество архитектуры MoE состоит в разделении двух величин: общей емкости модели и числа операций на один вход. Модель может содержать много экспертов и, следовательно, иметь большую параметрическую емкость, но конкретный токен обрабатывается только выбранным подмножеством экспертов. Это делает MoE привлекательным направлением развития больших языковых моделей: оно позволяет увеличивать масштаб модели без такого же резкого роста вычислительной стоимости инференса.

Эта идея получила развитие в крупных трансформерных моделях. GShard использует условное вычисление и автоматическое разбиение вычислений для масштабирования MoE трансформеров до сотен миллиардов параметров~\cite{lepikhin2020gshard}. Switch Transformer упрощает маршрутизацию, выбирая одного эксперта для каждого токена, и показывает возможность масштабирования до триллионного числа параметров при контролируемой вычислительной стоимости~\cite{fedus2021switch}. GLaM также использует разреженно активируемую архитектуру MoE и показывает, что она может увеличивать емкость языковой модели при меньшей стоимости обучения по сравнению с плотными вариантами сопоставимого качества~\cite{du2021glam}. DeepSpeed-MoE рассматривает уже не только обучение, но и инференс MoE моделей, подчеркивая, что для практического применения нужны специальные системные решения~\cite{rajbhandari2022deepspeedmoe}.

При этом MoE не устраняет системные сложности, а переносит их в другую часть системы. Вместо одной плотной операции возникает маршрутизация токенов, распределение нагрузки между экспертами, отправка активаций к выбранным экспертам и сбор результатов. Если эксперты размещены на разных GPU или узлах, MoE слой становится не только вычислительной, но и коммуникационной задачей.

\subsection{Методы распределенного инференса}

Распределенный инференс нужен в двух основных случаях. Первый случай --- модель или ее рабочее состояние не помещаются в память одного GPU. Второй случай --- одного устройства недостаточно по скорости, даже если памяти хватает. В обоих случаях модель, входные данные или поток запросов распределяются между несколькими устройствами, а система должна согласовать вычисления и обмен промежуточными результатами.

Параллелизм по данным применяется, когда нужно обслуживать больше независимых запросов без изменения устройства одной копии модели. В этом случае несколько копий модели запускаются на разных GPU, а пользовательские запросы распределяются между этими копиями. Такой подход хорошо работает, когда модель помещается в память одного устройства, однако не помогает, если модель слишком большая для одной видеокарты.

Тензорный параллелизм применяется, когда отдельные матричные операции внутри слоя разбиваются между несколькими GPU. Например, части весов могут храниться на разных устройствах, каждое устройство считает свою часть результата, после чего результаты объединяются с помощью коллективной операции (All-Reduce). Этот подход является стандартным способов распределенного инференса крупных моделей. В статье Megatron-LM такая схема описана для разбиения на уровне слоя модели на множестве GPU~\cite{shoeybi2019megatron}; та же идея используется и в системах инференса, где нужно уменьшить объем памяти и вычислений на одном устройстве.

Конвейерный параллелизм размещает группы слоев на разных устройствах. Когда одна группа слоев обрабатывает вход, она передает активации следующей группе, и так далее. Такой подход снижает требования к памяти одного GPU, но добавляет зависимость между стадиями. Для обучения конвейер можно заполнять микробатчами, а при инференсе отдельных запросов он потенциально может увеличивать задержку, так как запрос должен пройти последовательную цепочку групп слоев.

Экспертный параллелизм характерен именно для MoE моделей. В нем разные эксперты размещаются на разных GPU или узлах, а токены после маршрутизации отправляются к тем экспертам, которые были выбраны роутером. В отличие от тензорного параллелизма, где коммуникационная схема часто регулярна и связана со структурой слоя, в экспертном параллелизме трафик зависит от входных данных. Это делает нагрузку менее предсказуемой: один эксперт может получить много токенов, другой --- мало или не получить ни одного.

На практике эти методы часто комбинируются. Например, плотные части модели могут использовать тензорный или конвейерный параллелизм, а экспертные слои --- экспертный параллелизм. Такой гибридный подход нужен, потому что разные части модели имеют разные ограничения: где-то не хватает памяти, где-то требуется ускорить матричные операции, а где-то главной проблемой становится обмен активациями между экспертами.

\subsection{Коллективные коммуникации и ограничения распределенного инференса}

Распределенный инференс невозможен без коммуникационных примитивов, то есть операций обмена данными между процессами-участниками. В литературе и практических системах широко используются коллективные операции, например Broadcast, Gather, AllReduce, AllGather, ReduceScatter и All-to-All. На NVIDIA GPU такие операции обычно реализуются с опорой на библиотеки из экосистемы NVIDIA. NCCL предоставляет коллективные и point-to-point операции для GPU~\cite{nvidiaNCCLDocs}. NVSHMEM реализует модель OpenSHMEM для кластеров NVIDIA GPU, поддерживает односторонний доступ к памяти, синхронизацию и коллективные операции из CPU, CUDA streams и непосредственно из CUDA kernels~\cite{nvidiaNvshmem}. Эти библиотеки важны не как самостоятельные архитектуры инференса, а как практическая основа, на которой строятся как обобщенные коллективные операции, так и более специализированные MoE коммуникации, включая dispatch/combine в библиотеках экспертного параллелизма. Для тензорного параллелизма типичны AllReduce, AllGather и ReduceScatter, а для экспертного параллелизма особенно важен All-to-All: каждый участник может отправить часть токенов каждому другому участнику, владеющему нужными экспертами.

Коллективные операции удобны тем, что дают высокопроизводительную абстракцию над сложной топологией GPU и сети. Разработчику не нужно вручную описывать каждую отдельную передачу: достаточно задать операцию для группы участников, а библиотека выбирает алгоритмы и транспорт. Для синхронного выполнения на фиксированном наборе устройств это естественная модель: все участники входят в одну коммуникационную группу, вызывают одну и ту же коллективную операцию и продолжают работу после ее завершения.

Однако именно эта связанность становится ограничением для задач, где решение о выборе исполнителя должно приниматься во время работы. Коллективная операция предполагает согласованное участие процессов из заданной группы. Если один участник задержался или стал недоступен, вся операция может задержаться или завершиться ошибкой. Современные библиотеки добавляют средства обработки ошибок, завершения коммуникационных групп и создания новых групп, но это не превращает коллективную операцию в механизм динамической балансировки нагрузки. Изменение состава участников требует явного управления группами и согласования состояния между оставшимися процессами~\cite{nvidiaNCCLDocs}.

Для инференса MoE моделей это особенно важно. В классической схеме экспертного параллелизма все GPU с экспертами участвуют в операции All-to-All. Такая схема эффективна, когда все участники исправны, а нагрузка на отдельные модули постоянна и распределена равномерно. Однако маршрутизация токенов в MoE зависит от входных данных, поэтому разные эксперты могут получать существенно разное число токенов. Показательно, что для этой проблемы существуют отдельные системные решения, например EPLB (Expert Parallelism Load Balancer). EPLB балансирует экспертный параллелизм за счет выбора числа реплик экспертов и их размещения по GPU и узлам на основе измеренной нагрузки~\cite{deepseekEPLB}.

Однако такая балансировка остается привязанной к заранее выбранной раскладке экспертов по рангам. Распределение входных запросов может меняться во времени, очереди у исполнителей могут расти неравномерно, а часть экспертных блоков может временно становиться недоступной или перегруженной. В этих условиях статически выбранное размещение экспертов остается лишь приближением к текущей нагрузке. При такой неравномерности быстрые участники вынуждены ждать медленных, а если экспертный блок нужно временно исключить или добавить во время работы новую реплику, то фиксированная коллективная схема не удовлетворяет таким требованиям. Если же один и тот же эксперт размещен в нескольких экспертных блоках, клиент-серверная схема позволяет трактовать эти блоки как адресуемые реплики и выбирать менее загруженную доступную реплику во время обработки запроса.

\subsection{DeepEP и коллективы для MoE}

Для MoE существуют и более специализированные реализации коллективного подхода. Обычная All-to-All операция передает наборы тензорных фрагментов между всеми участниками, но MoE слой имеет более конкретную структуру: сначала выполняется dispatch, то есть отправка токенов к выбранным экспертам, затем вычисление экспертов, затем combine, то есть сбор результатов обратно в исходный порядок. Поэтому специализированная библиотека может оптимизировать не абстрактный All-to-All примитив, а именно пару dispatch/combine с учетом формы тензоров, числа экспертов, top-k маршрутизации, точности данных и используемой межузловой или внутрисерверной связи.

Одной из таких реализаций является DeepEP. Библиотека DeepEP сфокусирована на экспертном параллелизме и предоставляет высокопроизводительные GPU-ядра для MoE dispatch и combine. Для нее важны режимы высокой пропускной способности и низкой задержки, поддержка низкой точности, включая FP8, а также применимость как к обучению, так и к инференсу современных MoE моделей~\cite{deepseekDeepEP}. DeepEP использует NVSHMEM как одну из низкоуровневых основ для коммуникаций между участниками и остается примером специализированной реализации MoE коммуникаций поверх низкоуровневого слоя обмена данными.

В данной работе DeepEP важен как представитель коллективного подхода к MoE инференсу. Предлагаемый метод противопоставляет две архитектурно разные схемы: синхронную схему экспертного параллелизма на коллективах и клиент-серверную схему с независимыми экспертными блоками. Поэтому для экспериментального сравнения не требуется сравниваться со всеми возможными реализациями All-to-All. Достаточно выбрать сильную специализированную реализацию, которая уже оптимизирована именно под MoE dispatch/combine и может рассматриваться как практически готовый к промышленному использованию ориентир для коллективного подхода. В этой роли DeepEP подходит лучше, чем простой All-to-All через обобщенную коммуникационную библиотеку, поскольку он ближе к тому, как MoE коммуникации оптимизируются в реальных системах.

\subsection{Движки адресной передачи памяти: NIXL и MoonCake Transfer Engine}

Переход от коллективной схемы MoE к клиент-серверной архитектуре не устраняет необходимость передавать данные между GPU. Наоборот, эта необходимость становится явной частью проектируемой системы. В коллективном подходе операция All-to-All не только задает синхронизацию участников, но и берет на себя перемещение байтов: активации доставляются к GPU с нужными экспертами, а результаты возвращаются обратно. Если заменить такую операцию на обращение основной части модели к конкретному экспертному блоку, то система должна самостоятельно решить, как эффективно передать область GPU-памяти от одного исполнителя к другому и затем получить результат.

Обычного управляющего протокола для этого недостаточно. Он может передать команду, идентификатор запроса или метаданные, но сами активации и результаты экспертных вычислений являются буферами в памяти GPU. Их невыгодно переносить через лишние промежуточные копирования и синхронные операции на CPU: при каждом обращении к экспертному блоку это напрямую увеличивает задержку MoE слоя. Поэтому в предлагаемой архитектуре нужен отдельный слой адресной передачи памяти между GPU.

Для обмена данными между видеокартами существует несколько уровней программных средств. На нижнем уровне находятся физические каналы и сетевые технологии: PCIe, NVLink~\cite{nvidiaNvlink} и InfiniBand. Выше располагаются транспортные библиотеки, например UCX, которые скрывают детали конкретного канала и позволяют работать с памятью CPU и GPU через общий интерфейс~\cite{openucxFeatures}. Еще выше находятся специализированные движки передачи данных. В контексте данной работы они нужны не как самостоятельная система инференса, а как замена той части коллективной операции, которая обеспечивала собственно перемещение буферов между GPU.

NIXL (NVIDIA Inference Xfer Library) является одной из реализаций такого слоя. Он используется в экосистеме NVIDIA Dynamo и ориентирован на перемещение данных в распределенном инференсе моделей~\cite{nvidiaNixl}. Dynamo --- открытый фреймворк для распределенного инференса LLM; NVIDIA представляет его как промышленную основу для масштабируемого инференса и указывает интеграции с TensorRT-LLM, SGLang и vLLM~\cite{nvidiaDynamoProduction}. В Dynamo NIXL применяется как механизм передачи KV кэша между исполнителями в disaggregated serving~\cite{nvidiaDynamoDisaggregatedServing}. Для предлагаемого MoE инференса важен тот же принцип: участники заранее обмениваются метаданными о доступных областях памяти, регистрируют буферы и затем выполняют адресные операции чтения или записи между конкретными сторонами. Такая модель позволяет основной части модели отправлять активации выбранному экспертному блоку, не вовлекая в обмен всю коллективную группу.

MoonCake Transfer Engine --- другая реализация того же класса задач. Она появилась в контексте систем обслуживания больших языковых моделей и используется как слой передачи данных в проекте MoonCake. Для нее важна поддержка разных протоколов: TCP, RDMA, NVMe-oF, а также транспортов на основе NVLink для связи внутри сервера и между связанными узлами~\cite{mooncakeSupportedProtocols}. В отличие от обычного обмена сообщениями, Transfer Engine предоставляет операции чтения и записи фрагментов DRAM/VRAM через выбранный транспорт~\cite{mooncakeTransferEngine}. Поэтому его можно использовать как основу для передачи памяти между основной частью модели и экспертными блоками, а не только управляющих сообщений.

В данной работе NIXL и MoonCake Transfer Engine рассматриваются как две конкретные реализации слоя адресной передачи памяти между GPU. Их использование позволяет отделить клиент-серверную архитектуру MoE слоя от конкретного механизма перемещения данных. Если метод работает только с одним движком, трудно понять, является ли результат следствием архитектуры или частной реализации передачи.

\subsection*{Выводы по главе}
\addcontentsline{toc}{subsection}{Выводы по главе}

В главе показано, что для инференса больших языковых моделей особенно важны задержка, пропускная способность, работа с памятью GPU и способность обслуживать поток запросов без падения всей системы. По мере роста моделей распределенный инференс становится не оптимизацией, а необходимым условием практического применения.

Архитектура MoE является одним из наиболее перспективных направлений развития больших языковых моделей. Но эта выгода достигается ценой усложнения маршрутизации и коммуникаций.

Традиционные коллективные коммуникации хорошо подходят для синхронного выполнения на фиксированной группе участников, но плохо сочетаются с требованием динамически выбирать менее загруженную реплику эксперта, исключать недоступные экспертные блоки и добавлять новых исполнителей во время работы. Поэтому для предлагаемого метода важен слой адресной передачи GPU-памяти: после отказа от All-to-All именно он должен обеспечить эффективное перемещение активаций и результатов между основной частью модели и выбранными экспертными блоками. NIXL и MoonCake Transfer Engine дают такую основу и позволяют строить архитектуру, в которой экспертные блоки являются более независимыми, наблюдаемыми и управляемыми компонентами распределенного инференса.
